\section{The original column values and the two global slope bounds}
\label{sec:global-slope}

\status{This chapter supplies the three GL4 entries retained in RSL13:
the actual column-value dictionary, the lower slope bound, and the
maximum slope bound with coefficient $4$.  The finite arguments below
are on the original lattices and metrics.  The one additional published
input is Banaszczyk's Euclidean transference theorem, with its hypotheses
checked in GLS8.  The common compatible family and its original second
layer remain inputs.}

\paragraph{Fixed objects and normalization.}
Fix the original record $D$ and a word $u$.  Write
\[
\begin{gathered}
 E=D.\mathrm{host}.E,\quad K=D.\mathrm{what}_2.\mathrm{LL},\quad
 k=N(0)=D.R.k,\quad d_E=[E:\mathbb Q]=2f,\\
 T=(D.\mathrm{mem}\,u).A.T,\qquad
 v=D.\mathrm{host}.v=\operatorname{rd}(E)/2.
\end{gathered}
 \tag{GLS.1}
\]
The class set is $\mathcal C=(D.\mathrm{mem}\,u).C.\mathrm{cls}(0)$.
Put $S=D.\mathrm{sp}.X(0)$, $L_c=D.\mathrm{clsLat}(u,0,c)$, and
$A=(D.\mathrm{fm}\,u).\mathrm{lat}(0)$.  All these objects, including
$A$'s integral lattice and all its infinite-place metrics, are fixed.
For the original dominant integral weight $\lambda_\tau$ set
\[
 s_\tau=\lambda_\tau(k-1),\quad
 \nu_{\tau,j}=\lambda_\tau(j)-s_\tau,\quad
 d_\tau=\sum_{j=0}^{k-1}\nu_{\tau,j},\quad
 a_\tau=\sum_{j=0}^{k-1}\lambda_\tau(j),\quad
 h_\tau=\nu_{\tau,0}.
 \tag{GLS.2}
\]
Thus $a_\tau,d_\tau,h_\tau$ are respectively the source's
\texttt{colAbs}, \texttt{colAbsP}, and \texttt{firstColumnSpread}.
The first of these is a signed sum.  The last two are nonnegative.
We use the source's squared Hermitian slope: for a rank-$r$ lattice
$M$ over a totally complex field $K$, a rational basis $e$, local basis
matrices $\Gamma_\mathfrak p$, and positive Gram matrices $H_\sigma$,
\[
 r\mu(M)=-\frac{2}{[K:\mathbb Q]}
       \sum_\mathfrak p\log N\mathfrak p\,
                      v_\mathfrak p(\det\Gamma_\mathfrak p)
       -\frac1{n_\infty(K)}\sum_\sigma\log\det H_\sigma.
 \tag{GLS.3}
\]
Here $n_\infty(K)=[K:\mathbb Q]/2$.  Only finitely many finite-place
terms are nonzero.  Formula (GLS.3) is the determinant/index formula
already used in MT2, with the same factor $2$.

\subsection{The same column-value record, at every place}

% BEGIN GLS1
\paragraph{GLS1: the original class lattices and their scalar extensions.}
\label{claim:GLS1}
Suppose an original \texttt{ColsTiedSym D} record $\mathcal F$ is
supplied.  Its actual definite space $S$ is standard positive at every
embedding.  Its level-zero class set is nonempty: the identity adelic
class has a preimage under the stored class bijection.  The original
lattice $L_c$ is integral, and
\[
 L_c\subseteq L_c^\vee,\qquad
 L_c^\vee/L_c\simeq\mathcal O_E/T\mathcal O_E,
 \qquad [L_c^\vee:L_c]=T^{2f}.                 \tag{GLS.4}
\]
Here $L_c^\vee=\{y:h_S(L_c,y)\subseteq\mathcal O_E\}$ is exactly
the source's \texttt{dualSub}.  To check the first assertion on the
original lattice, the stored local representative identifies its
completion with the adelic translate of the original integral base
lattice.  A unitary translate preserves inclusion in the Hermitian
dual.  Intersecting these local inclusions with $S$ gives the first
inclusion in (GLS.4).  The stored \texttt{atT.disc} identifies the
actual quotient with the displayed residue module, as in RMP2; its
cardinality is $T^{[E:\mathbb Q]}$.  This proves the whole literal
\texttt{LatT} condition, including its index.

Choose the source embedding
\[
 \rho_\tau=\mathrm{embE}\circ\mathrm{conj}_E\circ
                   \sigma_{\mathrm{eTau}(\tau)}:E\longrightarrow K.
 \tag{GLS.5}
\]
For every $\sigma:K\hookrightarrow\mathbb C$ it satisfies
$\sigma\rho_\tau=\overline{\mathrm{embS}(D,\sigma,\tau)}$.
Let $G$ be the Hermitian Gram matrix of the original frame of $S$,
reindexed by $\mathrm{Fin}(k)$.  Define $Q_{c,\tau}$ on $K^k$ by
\[
 \operatorname{bc}_\tau(z)=\rho_\tau(\operatorname{coord}_S z),
 \qquad Q_{c,\tau}.L=
       \mathcal O_K\operatorname{-span}\operatorname{bc}_\tau(L_c),
 \qquad H_{c,\tau,\sigma}=\sigma\rho_\tau(G).
 \tag{GLS.6}
\]
Choose an invertible complex matrix $X_{c,\tau,\sigma}$ with
$X_{c,\tau,\sigma}^{-*}X_{c,\tau,\sigma}^{-1}
=H_{c,\tau,\sigma}$.  The positive Hermitian square root gives one.
The original full lattice $L_c$ has finitely many generators; their
images generate the lattice in (GLS.6).  Its $K$-span is $K^k$ because
the original rational frame is in the $E$-span of $L_c$.  Thus each $Q$
is good.  The coordinate map is additive and $\rho_\tau$-semilinear,
and at every infinite place
\[
 \|\operatorname{bc}_\tau(z)\|_{Q,\sigma}^{2}
       =\operatorname{Re}\sigma\rho_\tau(h_S(z,z)).
 \tag{GLS.7}
\]
This constructs all fields of an actual \texttt{ColBase D u} with
\texttt{rho = rhoC D}.  In particular it uses the original $L_c$,
and (GLS.7) holds at every coefficient embedding.
% END GLS1

% BEGIN GLS2
\paragraph{GLS2: the coefficient map and the exact local criterion.}
\label{claim:GLS2}
Assume additionally \texttt{adm3 D}; its assertion at $u$ says that
all original level-zero stabilizers are trivial.  Project
$\mathcal F.\mathrm{forms}(u)$ to its original \texttt{Forms0} record.
The rational-value dictionary and polynomial injectivity give a
$K$-linear map
\[
 \operatorname{co}_0:A.V\longrightarrow
       \bigoplus_{c\in\mathcal C}W_K,\qquad
 W_K=\bigotimes_\tau \Delta_{\nu_\tau}\otimes_{\mathbb Z}K.
 \tag{GLS.8}
\]
The polynomial variables for different $\tau$ are disjoint.  The
integral lattice $\Delta_\nu$ is the divided-power Weyl lattice
generated by its actual highest vector; it is the literal
\texttt{delSq}, rather than an arbitrary full polynomial lattice.
The value of $\operatorname{co}_0(x)_c$ on the original frame chart is
the original rational form $x$ on class $c$.  Vanishing of all these
polynomials implies vanishing of the form.  Conversely, any tuple in
$W_K$ gives an allowed function on each class: the stabilizer
equivariance condition is empty, and the stored rational-span equality
puts that tuple in $A.V$.  Polynomial uniqueness then gives (GLS.8)
as an isomorphism.  This is also the original rational surjectivity
used in GLE4 and RSL9.

Use exactly the source action
\[
 \mathcal R(g)P(Z)=\left(\prod_\tau\det(g_\tau)^{s_\tau}\right)
                              P((g_\tau^{t}Z_\tau)_\tau).
 \tag{GLS.9}
\]
Set $g_{c,\tau}=\rho_\tau(G)^{-1}$ for the coordinates (GLS.6), and
$\operatorname{co}(x)_c=\mathcal R(g_c)\operatorname{co}_0(x)_c$.
This is still an isomorphism onto the same polynomial space.  For a
general coordinate basis of $Q$ the formula is
$g_{c,\tau}=M_{c,\tau}\rho_\tau(G)^{-1}$, where $M$ is the matrix
of the original map $\operatorname{bc}$; this is precisely
\texttt{gC} in the source.

Fix an arbitrary coefficient place $w$, not just one over $T$.
Choose an $E$-rational basis matrix $Y_c$ for $L_c$ simultaneously at
the finitely many places over the rational residue characteristic of
$w$.  Such a basis exists by local freeness over the Dedekind domain
and the Chinese remainder theorem: prescribe a basis modulo each of
these primes, choose the columns in $L_c$, and use Nakayama at each
prime.  Put $\gamma_{c,\tau}=\rho_\tau(Y_c)$ in (GLS.6), or
$M_{c,\tau}\rho_\tau(Y_c)$ in the general coordinates.  These are
local bases of the exact scalar-extension lattice $Q_{c,\tau}.L$.
The corresponding original Hermitian dual basis has matrix
$(Y_c^*G)^{-1}$.  At the distinguished coefficient embedding,
\[
 \left(\mathrm{embLL}
             (\gamma_{c,\tau}^{-1}g_{c,\tau})\right)^t
       =\left(\mathrm{embT}_\tau(Y_c^*G)\right)^{-1}.
 \tag{GLS.10}
\]
Indeed $\gamma^{-1}g=\rho_\tau(GY_c)^{-1}$, and (GLS.5)
turns the transpose of its image into the inverse of $Y_c^*G$.
Substitute (GLS.10) in the original \texttt{fnOf} formula.
Its determinant power is exactly $s_\tau$, so both its polynomial
substitution and its scalar determinant factor agree with (GLS.9).
The original local integral-value criterion of GLE4 consequently says
\[
 x\in A.L_w\quad\Longleftrightarrow\quad
 \mathcal R(\gamma_c^{-1})\operatorname{co}(x)_c
       \in\bigotimes_\tau\Delta_{\nu_\tau,\mathcal O_w}
       \quad\hbox{for every }c.                  \tag{GLS.11}
\]
The right side is the source's \texttt{IntSpanAt}.  GLE3 proves its
invariance under every integral invertible change of basis, including
the signed determinant factor.  Any other local basis of the same
$Q_{c,\tau}.L$ differs from $\gamma_{c,\tau}$ by such a change.
Thus (GLS.11) holds for \emph{every} family of local bases.  This
checks the full quantifier in \texttt{ColVal.fin}.
% END GLS2

% BEGIN GLS3
\paragraph{GLS3: the original all-place norm, and a literal ColVal.}
\label{claim:GLS3}
At an infinite place $\sigma$ put
\[
 Z_{c,\tau}=X_{c,\tau,\sigma}^{-1}\,
                         \sigma(g_{c,\tau}).
 \tag{GLS.12}
\]
In general coordinates insert $M$ as in GLS2.  The identity
$(X^{-1}\sigma M)^*(X^{-1}\sigma M)=\sigma\rho_\tau(G)$,
together with (GLS.5), gives
\[
 (Z_{c,\tau}^{t})^*\,
      \mathrm{embS}(D,\sigma,\tau)(G)\,Z_{c,\tau}^{t}=1.
 \tag{GLS.13}
\]
This follows by multiplying out $Z=Y(\sigma\rho_\tau G)^{-1}$
with $Y^*Y=\sigma\rho_\tau G$ and transposing.  Thus $Z_c^t$ is
an orthonormal frame for the original definite form space at this
embedding.

Every class weight in the original norm formula is
$1/|\mathrm{stab}(0,c)|=1$.  To allow the orthonormal frame in
(GLS.13) to depend on $c$, use (GLS.8) to take the rational form
supported on that single class.  Its original \texttt{norm0} formula
shows that this class's contribution is independent of the chosen
orthonormal frame.  Apply this separately to each class and add.
The determinant term in \texttt{fnOf} is retained throughout;
$\det(Z^t)=\det Z$ and (GLS.9) includes its exact signed power.
One obtains
\[
 \|x\|_{A,\sigma}^{2}
 =\sum_{c\in\mathcal C}
  \frac{\left\|
    \mathcal R(X_{c,\sigma}^{-1})
                \sigma(\operatorname{co}(x)_c)\right\|_F^2}
       {\prod_\tau\|H_{\nu_\tau}\|_F^2}.        \tag{GLS.14}
\]
Here the Fock squared norm is $\|P\|_F^2=\sum_\alpha
|P_\alpha|^2\alpha!$, with the tensor convention on disjoint
blocks.  The equality follows from
$\mathcal R(X^{-1})\sigma(\operatorname{co}x)
=\mathcal R(Z)\sigma(\operatorname{co}_0x)$.

GLS1 supplies the nonempty class set, the exact \texttt{LatT}, and
all scalar-extension fields; GLS2 supplies the rational map, its
range, and its local criterion; (GLS.14) supplies the last metric
field.  These are every field of the literal \texttt{ColVal D u}.
We have therefore proved, pointwise on the original record,
\[
 \texttt{ColsTiedSym}(D)\ \wedge\ \texttt{adm3}(D)
       \Longrightarrow\ \forall u,\
                          \texttt{Nonempty} (\texttt{ColVal}(D,u)).
 \tag{GLS.15}
\]
This is the exact \texttt{ColumnValuationFromTwistCompatibility}
contract.  No column values, class representatives, or integral
lattice were reselected independently of the supplied forms.
% END GLS3

\subsection{Two finite polynomial estimates with their normalization}

% BEGIN GLS4
\paragraph{GLS4: the highest-vector norm and the hook bound.}
\label{claim:GLS4}
For a partition $\nu_1\ge\cdots\ge\nu_k\ge0$ use one-based indices
in this paragraph.  Put
$H_\nu=\prod_{j=1}^k\det(Z_{[j],[j]})^{\nu_j-\nu_{j+1}}$,
where $\nu_{k+1}=0$, and let $N_\nu=F(H_\nu,H_\nu)$ for the
bilinear rational Fock form.  Then
\[
 N_\nu=
 \frac{\prod_{i=1}^k(\nu_i+k-i)!}
      {\prod_{1\le i<j\le k}(\nu_i-\nu_j+j-i)}.
 \tag{GLS.16}
\]
We give the finite differential calculation behind the formula.
Multiplication by $z_{as}$ is Fock-adjoint to $\partial_{as}$, as
is seen on two monomials.  Set $E_{ab}=\sum_s z_{as}\partial_{bs}$.
The following normal-order identity holds, with factors in increasing
column order:
\[
 \sum_{\pi\in\mathfrak S_k}\operatorname{sgn}(\pi)
 \prod_{j=1}^k\bigl(E_{\pi(j),j}+(k-j)\delta_{\pi(j),j}\bigr)
           =\det Z\,\det\partial.              \tag{GLS.17}
\]
For completeness, this can be verified without a representation
theorem.  For lists $w,g$ of length $r$ let
\[
 P_r(w,g)=\sum_{s_1,\ldots,s_r}
 z_{w_1s_1}\cdots z_{w_rs_r}
 \partial_{g_1s_1}\cdots\partial_{g_rs_r},\qquad
 L_r(w,g)=\prod_{j=1}^r
       (E_{w_jg_j}+(r-j)\delta_{w_jg_j}).
\]
The relation $\partial_{bs}z_{at}=z_{at}\partial_{bs}
+\delta_{ab}\delta_{st}$ gives
\[
 P_r=E_{w_1g_1}P_{r-1}(w',g')
   -\sum_{j=2}^r\delta_{w_jg_1}
                   P_{r-1}(w'[j\leftarrow w_1],g').
 \tag{GLS.18}
\]
Subtract this recurrence from that for $L_r$.  Besides
$(E_{w_1g_1}+(r-1)\delta_{w_1g_1})(L_{r-1}-P_{r-1})$, the
difference has $(r-1)\delta_{w_1g_1}P_{r-1}$ and the displayed
sum with positive sign.  Antisymmetrize in $w$.  The first term
vanishes by induction.  For each $j$, the transposition of positions
$1,j$ changes the sign and identifies the $j$th term with
$\delta_{w_1g_1}P_{r-1}$; the remaining $r-1$ terms cancel it.
Thus the antisymmetrizations of $L_r$ and $P_r$ agree.  When
$w=g=(1,\ldots,k)$, terms with repeated $s_j$ vanish in the
determinant, and the remaining permutations give
$\det Z\det\partial$.  This proves (GLS.17).

On $H_{\nu+\mathbf1}$, $E_{ii}$ acts by $\nu_i+1$, and
$E_{ij}$ kills it when $i<j$; both statements follow directly from
the leading-minor product.  In a nonidentity permutation term of
(GLS.17), take the largest moved column $j$.  Then $\pi(j)<j$;
all factors to its right are diagonal, so that term kills the
highest vector.  Consequently (GLS.17), cancellation of the nonzero
polynomial $\det Z$, and Fock adjointness give
\[
 N_{\nu+\mathbf1}
          =\prod_{i=1}^k(\nu_i+k-i+1)\,N_\nu.
 \tag{GLS.19}
\]
If $\nu_k=0$, $H_\nu$ uses only the first $k-1$ rows and columns;
its Fock norm is unchanged on deleting the last row and column.
These two recurrences prove (GLS.16), by induction on $k$ and then
on $\nu_k$: the right side of (GLS.16) satisfies exactly the same
recurrences and is $1$ in dimension zero.

The coefficient of $\prod_i z_{ii}^{\nu_i}$ in $H_\nu$ is $1$.
Its Fock contribution alone is $\prod_i\nu_i!$.  Thus, with the
source normalization $c_\nu=N_\nu/\prod_i\nu_i!$,
\[
 1\le c_\nu\le(\nu_1+k)^{k(k-1)/2},\qquad
 \log c_\nu\le k^2\log(\nu_1+k).             \tag{GLS.20}
\]
For the upper bound, each positive integer factor in the denominator
of (GLS.16) is at least one, and
$(\nu_i+k-i)!/\nu_i!\le(\nu_1+k)^{k-i}$.
The last logarithmic inequality uses $\nu_1+k\ge1$.
This proves the literal \texttt{HookBound}, with its exact Fock
normalization.
% END GLS4

% BEGIN GLS5
\paragraph{GLS5: the normalized integral Weyl slope.}
\label{claim:GLS5}
If $k=0$, the only polynomial is a constant:
$\Delta=\nabla=\mathbb Z\cdot1$, its normalized Gram matrix is
$(1)$, $c_\nu=1$, and both the slope and the required lower bound
are zero.  Thus suppose $k\ge1$, put $d=\sum_j\nu_j$, and let
$r=\dim_\mathbb Q W_\nu>0$.
Let $\nabla_\nu$ be the integral coefficient lattice in the same
polynomial space.  IW5 gives, for the normalized partition, the
exact duality
\[
 \Delta_\nu^{\#}=\nabla_\nu,\qquad
 B_\nu=F/N_\nu.                              \tag{GLS.21}
\]
Here is the extension to every partition, including $\nu_k>0$.
Put $s=\nu_k$ and $\nu^0=\nu-s\mathbf1$. Multiplication by
$\det(Z)^s$ identifies the two rational highest-weight spaces,
and it maps $\Delta_{\nu^0}$ onto $\Delta_\nu$: every
off-diagonal root operator annihilates the determinant factor,
so it commutes with this multiplication, including its divided
powers.  It also maps $\nabla_{\nu^0}$ onto $\nabla_\nu$.
Indeed the rational spaces already match, and the polynomial
$\det(Z)^s$ is primitive; the DVR Gauss lemma shows that a rational
polynomial whose product with this factor has integral coefficients
itself has integral coefficients.  This also transports the
integral basis of CB2.

The pullback of $B_\nu$ has the same transpose-adjoint root
operators as $B_{\nu^0}$ and gives the highest vector norm $1$.
These properties determine the form: move each divided root word
in the first entry into the second by adjointness, and then take
the highest-weight component, which is a scalar multiple of the
highest vector.  Root words span the rational space, so all
pairings are determined.  Thus the pullback is exactly $B_{\nu^0}$,
and the two lattice identities transfer to the arbitrary partition.
This proves (GLS.21) in its full stated scope; it does not assume
that unnormalized Fock norms are unchanged by determinant
multiplication.

IW5 proves the normalized case at every DVR.  It holds over $\mathbb Z$ because
membership in a full rational lattice is equivalent to membership
at every prime.  In integral bases, the matrix of the pairing between
the two lattices is unimodular.  It follows that their positive Gram
determinants are reciprocal.

Here is the requisite bound on the coefficient lattice's determinant.
Every monomial $Z^\alpha$ occurring in it has column sums $\nu_j$.
The multinomial theorem gives
\[
 \frac{\prod_j\nu_j!}{\alpha!}\le k^d,
 \qquad \frac{\alpha!}{N_\nu}\ge
                        (k^d c_\nu)^{-1}=q_0.\tag{GLS.22}
\]
Choose a finite set of monomials containing all supports of an
integral basis of $\nabla_\nu$, and let $V$ be its integer
coefficient matrix.  Its normalized Gram matrix is
$V^t\operatorname{diag}(\alpha!/N_\nu)V$.
Cauchy--Binet expresses its determinant as the sum of squared
integer $r$-minors of $V$, weighted by products of $r$ diagonal
entries.  Some minor is nonzero, since the basis is independent.
Every weight is at least $q_0^r$, so the determinant is at least
$q_0^r$.  By (GLS.21), the Gram determinant on $\Delta_\nu$ is at
most $q_0^{-r}$.  Consequently
\[
 \mu_{\rm std}(\nu)
 :=-\frac1r\log\det\operatorname{Gram}_{B_\nu}(\Delta_\nu)
 \ \ge\ -d\log k-\log c_\nu.                 \tag{GLS.23}
\]
This is the source's \texttt{StdSlopeLow}.  It also covers $d=0$:
the polynomial lattice is the one-dimensional lattice generated by
$1$, and both sides are zero.  All determinants here use the literal
divided-power lattice and the normalization of (GLS.14).
% END GLS5

\subsection{The lower slope on an arbitrary original ColVal}

% BEGIN GLS6
\paragraph{GLS6: exact slope of every column lattice in a supplied ColVal.}
\label{claim:GLS6}
Now take \emph{any} $\mathrm{CV}:\texttt{ColVal D u}$, without
assuming it was constructed by GLS1.  Its semilinear map extends to
a $K$-linear map $K\otimes_{E,\rho_\tau}S\to Q_{c,\tau}.V$.
The image contains the full lattice \texttt{CV.bc\_L}, so it is
surjective; both dimensions are $k$, so it is an isomorphism.
The metric condition \texttt{bc\_norm} identifies all diagonal
Hermitian values.  Polarization identifies the entire form: use
$z+z'$ for its real part and $z+az'$ for an $a\in E$ with nonreal
image to recover its imaginary part.  Such an $a$ exists because
$E$ is CM.  Every embedding of $E$ extends to $K$; hence positivity
of the good lattice $Q$ also proves positivity of $S$ at every
embedding.

The normalized slope is preserved by this scalar extension.  To
see every normalization, in a local basis of a full lattice a
determinant valuation at $\mathfrak p$ is multiplied by
$e(\mathfrak P/\mathfrak p)$ upstairs.  In (GLS.3) the sum
$\sum_{\mathfrak P\mid\mathfrak p}ef=[K:\rho_\tau(E)]$
cancels the increase of the field degree.  At infinity, each
complex place of $E$ has exactly $[K:E]$ places above it; a conjugate
choice has the same positive Gram determinant.  Thus the normalized
average in (GLS.3) is unchanged.  This argument applies to the
literal span \texttt{bc\_L}; no finite index is introduced.

For a positive Hermitian $E$-lattice $L$, the same determinant
formula gives $\mu(L^\vee)=-\mu(L)$: the dual ideal determinants
are inverse conjugates and the Gram matrices are inverse
transposes.  If $L\subseteq L^\vee$ has index $q$, the finite index
formula gives
\[
 \mu(L)=\mu(L^\vee)-\frac{2\log q}{d_E k},
 \qquad \mu(L)=-\frac{\log q}{d_E k}.          \tag{GLS.24}
\]
The \texttt{latT} field of this very $\mathrm{CV}$ has
$q=T^{2f}$ and $d_E=2f$.  Scalar extension therefore gives
\[
             \boxed{\ \mu(Q_{c,\tau})=-\frac{\log T}{k}\ }
                  \qquad(c\in\mathcal C,\ \tau\in\mathrm{Tau}).
 \tag{GLS.25}
\]
This proves the full \texttt{QSlope} contract for an arbitrary
\texttt{ColVal}, including its given embeddings and metrics.
% END GLS6

% BEGIN GLS7
\paragraph{GLS7: the exact polynomial determinant and the lower bound.}
\label{claim:GLS7}
Let $r_\tau=\operatorname{rank}_{\mathbb Z}\Delta_{\nu_\tau}$ and
$r=\prod_\tau r_\tau$.  Choose the tensor integral basis $Z$ of
$\bigotimes_\tau\Delta_{\nu_\tau}$ and let $G_0$ be its normalized
Fock Gram matrix.  Disjoint variables make this matrix a Kronecker
product, so
\[
 -\frac1r\log\det G_0=\sum_\tau\mu_{\rm std}(\nu_\tau).
 \tag{GLS.26}
\]
For the representation (GLS.9), its determinant satisfies
\[
 \det\mathcal R(g)^k=\prod_\tau\det(g_\tau)^{r a_\tau}.
 \tag{GLS.27}
\]
Indeed the integral polynomial lattice is graded by row weights.
Permutation matrices permute the rows, so the sum of the weights
of any one row over a homogeneous basis is $rd_\tau/k$.
The determinant twist adds $rs_\tau$, making the exponent
$ra_\tau/k$.  This proves (GLS.27) on diagonal matrices.
Elementary unipotents act unipotently and have determinant one;
diagonal matrices and these elementary matrices generate
$\mathrm{GL}_k$.  Thus the identity holds for every invertible
matrix, and its integer-power form avoids any root choices.

Use the rational basis $\operatorname{co}^{-1}(\mathbf1_cZ_i)$
of the original $A.V$.  At a finite place its integral basis is
$\operatorname{co}^{-1}(\mathbf1_c\mathcal R(\gamma_c)Z_i)$ by
\texttt{CV.fin}; at infinity its Gram blocks are
$R_{c,\sigma}^*G_0R_{c,\sigma}$ by \texttt{CV.inf}, with
$R_{c,\sigma}=\mathcal R(X_{c,\sigma}^{-1})$ in the basis $Z$.
Substitute these determinants and (GLS.27) into (GLS.3).  The finite
and infinite contributions belonging to a column combine into
$ra_\tau\mu(Q_{c,\tau})$.  Therefore
\[
 r|\mathcal C|\mu(A)
 =-|\mathcal C|\log\det G_0+
                 \sum_{c,\tau}ra_\tau\mu(Q_{c,\tau}).
 \tag{GLS.28}
\]
Here $r>0$ and $|\mathcal C|>0$; the latter is a field of
\texttt{ColVal}.  Divide (GLS.28), and apply GLS6 and (GLS.26):
\[
 \mu(A)=-\sum_\tau\frac{a_\tau}{k}\log T
                       +\sum_\tau\mu_{\rm std}(\nu_\tau).
 \tag{GLS.29}
\]
Using GLS4--GLS5 now proves exactly
\[
 \boxed{\quad
 -\sum_\tau\frac{a_\tau}{k}\log T
 -\sum_\tau\bigl(d_\tau\log k+k^2\log(h_\tau+k)\bigr)
                  \ \le\ \mu(A).
 \quad}                                                       \tag{GLS.30}
\]
This is \texttt{SlopeLow D u}, for every original $D,u$ and every
$\mathrm{CV}:\texttt{ColVal D u}$.  No size or eventuality
hypothesis was used.
% END GLS7

\subsection{The maximum slope with the exact coefficient four}

% BEGIN GLS8
\paragraph{GLS8: a published transference theorem on the exact trace lattice.}
\label{claim:GLS8}
For each original $L_c$ occurring in any \texttt{ColVal}, there are
$k$ $E$-linearly independent elements $y_1,\ldots,y_k\in L_c^\vee$
such that
\[
 \frac{\operatorname{Tr}_{E/\mathbb Q}h_S(y_a,y_a)}{2f}
                    \le k^2(2v)^2 \qquad(1\le a\le k).
 \tag{GLS.31}
\]
The published input is W.~Banaszczyk, \emph{New bounds in some
transference theorems in the geometry of numbers}, Math. Ann.
\textbf{296} (1993), 625--635, Theorem~2.1, p.~631,
\url{https://doi.org/10.1007/BF01445125}.
It states that a full Euclidean lattice $M\subset\mathbb R^n$
and its integer-pairing dual satisfy
$\lambda_i(M)\lambda_{n+1-i}(M^*)\le n$.
The case $n=1$ is equality.  This is the only result imported here
instead of being proved in this chapter.

Here is the complete matching of hypotheses and constants.
Use $M=L_c^\vee$, as a lattice of rank $n=d_E k$ in
$S\otimes_\mathbb Q\mathbb R$, with positive Euclidean form
\[
                  B(x,y)=\operatorname{Tr}_{E/\mathbb Q}h_S(x,y).
 \tag{GLS.32}
\]
It is real and symmetric because trace commutes with CM conjugation;
it is positive definite by GLS6.  The $\mathcal O_E$-dual of a full
projective lattice is full and finitely generated, so its Minkowski
image is discrete and spans this real vector space.  The integer
pairing dual in the theorem is exactly
\[
                  M^*=\mathfrak D_E^{-1}L_c.    \tag{GLS.33}
\]
Indeed $\operatorname{Tr}h(x,M)\subseteq\mathbb Z$ is equivalent
to $h(x,M)\subseteq\mathfrak D_E^{-1}$, since $M$ is
$\mathcal O_E$-stable.  Local duality in a basis, followed by
intersection over the primes, gives (GLS.33).

The original integrality $h(L_c,L_c)\subseteq\mathcal O_E$ and
$\overline{\mathfrak D_E}=\mathfrak D_E$ imply, for $0\ne z\in M^*$,
that $\beta=h_S(z,z)$ is totally positive and belongs to
$\mathfrak D_E^{-2}$.  Its principal fractional ideal multiplied by
$\mathfrak D_E^2$ is a nonzero integral ideal.  Hence
$|N_{E/\mathbb Q}\beta|\,|\operatorname{disc}E|^2\ge1$.
Arithmetic--geometric mean over the $d_E$ embeddings gives
\[
 B(z,z)\ge d_E/\operatorname{rd}(E)^2,
 \qquad \lambda_1(M^*)\ge\sqrt{d_E}/\operatorname{rd}(E).
 \tag{GLS.34}
\]
Apply the quoted theorem with $i=n$.  It gives $n$ independent
vectors in the original $M$ with squared norm at most
$d_Ek^2\operatorname{rd}(E)^2$; minima are attained because $M$
is discrete.  They span over $\mathbb Q$, hence over $E$.  Extract
an $E$-basis from them, preserving the bound, and divide it by
$d_E=2f$.  Since $\operatorname{rd}(E)=2v$, this is (GLS.31).
Thus the source's \texttt{ShortAt} witnesses are supplied in the
original dual lattice with the exact constant.
% END GLS8

% BEGIN GLS9
\paragraph{GLS9: coefficient ideals and the signed determinant correction.}
\label{claim:GLS9}
Fix a class $c$, its short dual basis from GLS8, and an arbitrary
$\mathrm{CV}$.  Work in the given rational coordinate basis of
$Q_{c,\tau}$.  The semilinear extension in GLS6 defines a unique
$K$-linear functional $\ell_{a,\tau}$ by
\[
 \ell_{a,\tau}(\operatorname{bc}_\tau z)
                         =\rho_\tau(h_S(y_a,z)).
 \tag{GLS.35}
\]
Let $g_\tau$ have these functionals as rows.  Nondegeneracy of $h_S$
and independence of the $y_a$ make $g_\tau$ invertible.  Every row
is integral on $Q_{c,\tau}.L$: this follows first on
$\operatorname{bc}_\tau(L_c)$ from $y_a\in L_c^\vee$ and
Hermitian symmetry, then on
its exact $\mathcal O_K$-span.  Define $\phi_c(x)$ to be the full
list of monomial coefficients of
\[
     P_c((g_\tau^{t}Z_\tau)_\tau),\qquad
                         P_c=\operatorname{co}(x)_c.
 \tag{GLS.36}
\]
All monomials have column multidegrees $\nu_\tau$; their finite
index set is exactly the source's \texttt{MonI}.  Invertible
substitution and the range in (GLS.8) make this coefficient map
injective on the subspace supported at $c$.

Let $M_\tau$ be the actual coordinate lattice of $Q_{c,\tau}.L$,
and let $\det M_\tau$ denote its determinant fractional ideal in
this basis.  Set
\[
 I_c=\prod_\tau(\det M_\tau)^{s_\tau},\qquad
 D_{c,\sigma}=\prod_\tau|\det X_{c,\tau,\sigma}|^{2s_\tau}>0.
 \tag{GLS.37}
\]
Negative $s_\tau$ are allowed.  At a finite place choose local
matrices $\gamma_\tau$ for $M_\tau$.  The coefficients of
$\mathcal R(\gamma^{-1})P_c$ are integral by
\texttt{CV.fin}, and $g_\tau\gamma_\tau$ has integral entries
by (GLS.35).  Substitution of $(g\gamma)^tZ$ shows that
\[
 \left(\prod_\tau\det\gamma_\tau^{-s_\tau}\right)
                   \phi_c(x)_\alpha\in\mathcal O_{K,w}
                \quad(x\in A.L\hbox{ supported at }c).
 \tag{GLS.38}
\]
Intersecting at all finite places gives
$\phi_c(x)_\alpha\in I_c$.  This is exactly the source's
\texttt{twId}/\texttt{twIdeal}: their definition intersects the
conditions (GLS.38) over every local basis.  Changes of local basis
have unit determinant and leave the same ideal.

Write $\mu_0(M_\tau)$ for the slope with the standard coordinate
metric.  Formula (GLS.3) yields both identities
\[
 \begin{split}
 -\frac{2}{[K:\mathbb Q]}\log N I_c
       &=\sum_\tau s_\tau k\mu_0(M_\tau),\\
 k\mu(Q_{c,\tau})
       &=k\mu_0(M_\tau)+\frac1{n_\infty(K)}
                            \sum_\sigma\log|\det X_{c,\tau,\sigma}|^2.
 \end{split}                                                   \tag{GLS.39}
\]
For the first equality the local exponent of $I_c$ is
$\sum_\tau s_\tau v_w(\det\gamma_\tau)$.
For the second, the coordinate Gram matrix is $X^{-*}X^{-1}$,
whose determinant is $|\det X|^{-2}$.
Combining these with GLS6 gives the exact cancellation
\[
 -\frac{2}{[K:\mathbb Q]}\log N I_c
 +\frac1{n_\infty(K)}\sum_\sigma\log D_{c,\sigma}
                       =-\sum_\tau s_\tau\log T.
 \tag{GLS.40}
\]
This computation keeps the original determinant shift and all
finite and infinite normalizations on the same $Q$.
% END GLS9

% BEGIN GLS10
\paragraph{GLS10: the class map's exact all-place norm estimate.}
\label{claim:GLS10}
Give the coefficient in monomial $\alpha$ the squared weight
\[
 q_\alpha=\frac{\alpha!}{\prod_\tau N_{\nu_\tau}},
 \qquad q_0=\left(k^{\sum_\tau d_\tau}
                                \prod_\tau c_{\nu_\tau}\right)^{-1}.
 \tag{GLS.41}
\]
The separate-column multinomial estimate of GLS5 proves
$q_\alpha\ge q_0>0$.  At an infinite place $\sigma$ let
\[
 G_{\tau,\sigma}=\sum_{a=1}^k
       \operatorname{Re}\sigma\rho_\tau(h_S(y_a,y_a))>0.
 \tag{GLS.42}
\]
The Riesz functional identity (GLS.35), in the orthonormal coordinates
given by $X_{c,\tau,\sigma}$, gives
\[
 \| (\sigma(g_\tau)X_{c,\tau,\sigma})^t\|_{\rm HS}^2
                                                   =G_{\tau,\sigma}.
 \tag{GLS.43}
\]
It follows by expanding $h_S$ in a basis and using the polarized
version of \texttt{bc\_norm}; the row of the dual functional is
the conjugate of the orthonormal coordinate vector of
$\operatorname{bc}_\tau(y_a)$.

For a polynomial of total degree $d_\tau$ in block $\tau$,
substitution by a matrix $U_\tau$ has Fock norm at most
$\|U_\tau\|_{\rm op}^{d_\tau}$ times the original norm.
One can check this directly by identifying a monomial of degree $d$
with its symmetric tensor: its Fock squared norm is $\alpha!$,
and the symmetric-tensor operator is the restriction of
$U^{\otimes d}$, whose norm is at most $\|U\|_{\rm op}^d$.
The normalization in this identification depends only on $d$ and
cancels.  Apply this separately to the blocks, and use
$\|U\|_{\rm op}\le\|U\|_{\rm HS}$.

In (GLS.36) write
$P=\mathcal R(X^{-1})\sigma(\operatorname{co}(x)_c)$.
Then (GLS.9) gives, including its signed determinant shift,
\[
 \sigma(\operatorname{co}(x)_c)((\sigma g)^tZ)
     =\left(\prod_\tau\det X_{c,\tau,\sigma}^{s_\tau}\right)
                     P(((\sigma g_\tau X_{c,\tau,\sigma})^tZ_\tau)_\tau).
\]
The $c$-summand of \texttt{CV.inf} is bounded by the whole norm.
Hence the coefficient map satisfies, on the original $A.V$,
\[
 \sum_\alpha q_\alpha|\sigma(\phi_c(x)_\alpha)|^2
 \le B_{c,\sigma}\|x\|_{A,\sigma}^2,
 \qquad B_{c,\sigma}=D_{c,\sigma}
                                \prod_\tau G_{\tau,\sigma}^{d_\tau}.
 \tag{GLS.44}
\]

All embeddings of $E$ extend equally often to $K$.  A conjugate
embedding leaves the real diagonal values in (GLS.42) unchanged.
Thus their average over $K$-places is the corresponding average
over $E$-places.  Since the full trace is twice the sum over the
complex places, (GLS.31) gives
\[
 \frac1{n_\infty(K)}\sum_\sigma G_{\tau,\sigma}
      =\sum_a\frac{\operatorname{Tr}h_S(y_a,y_a)}{2f}
      \le k^3(2v)^2.
\]
Concavity of $\log$ (or summing $\log t\le t-1$ after division
by the mean) gives
\[
 \frac1{n_\infty(K)}\sum_\sigma\log G_{\tau,\sigma}
                             \le\log(k^3(2v)^2).
 \tag{GLS.45}
\]
Equations (GLS.38), (GLS.40), (GLS.44), and (GLS.45) supply all
fields of the original \texttt{ClassMapT} and \texttt{ClassMap}.
% END GLS10

% BEGIN GLS11
\paragraph{GLS11: the maximum slope and the exact constant $4$.}
\label{claim:GLS11}
We first record the elementary slope comparison used for (GLS.44).
If a full lattice $N$ maps injectively to coefficient vectors in
$I^J$, with diagonal squared metric weights at least $q_0>0$, and
the squared norm of the map at $\sigma$ is at most $B_\sigma$, then
\[
 \mu_{\max}(N)\le-\log q_0
       -\frac{2\log NI}{[K:\mathbb Q]}
       +\frac1{n_\infty(K)}\sum_\sigma\log B_\sigma.
 \tag{GLS.46}
\]
Here is a determinant proof valid for every nonzero sublattice.
Write a rank-$r$ sublattice as $\bigoplus_j\mathfrak a_jz_j$ in
a pseudo-basis and choose a nonzero $r$-minor $\beta$ of its
coefficient vectors.  Integrality gives
$\beta\prod_j\mathfrak a_j\subseteq I^r$, whence
$|N\beta|\prod_jN\mathfrak a_j\ge(NI)^r$.
Cauchy--Binet bounds its coefficient Gram determinant below by
$q_0^r|\sigma\beta|^2$.  The product formula and (GLS.3) imply
that its coefficient-space slope is at most
$-\log q_0-2\log NI/[K:\mathbb Q]$.
Taking the $r$th exterior power in the norm bound gives the last
term in (GLS.46).  Taking the supremum over sublattices proves the
stated maximum-slope inequality.

The coordinate projection to a single class preserves $A.L$:
apply \texttt{CV.fin}, set the other class polynomials to zero,
and intersect the local integral lattices.  The projections sum to
the identity, and \texttt{CV.inf}, polarized, makes their ranges
orthogonal at every infinite place.  Thus $A.L$ is the actual
orthogonal direct sum of its class lattices $A_c$.
Its maximum slope is at most the maximum of theirs.  For example,
filter any sublattice by successive class projections; its successive
quotients inject into the individual $A_c$, with quotient norm at
least the projected norm.  Additivity of determinant degrees and
the same exterior-power comparison used for (GLS.46) give the
claimed bound.  The original class set is finite and nonempty.

Apply (GLS.46) to each $A_c$ with GLS9--GLS10.  Since
$-\log q_0=\sum_\tau(d_\tau\log k+\log c_{\nu_\tau})$,
the result before weakening the signed shift is
\[
 \mu_{\max}(A)\le
  \sum_\tau\left(d_\tau\log(4k^4v^2)
                    +\log c_{\nu_\tau}-s_\tau\log T\right).
 \tag{GLS.47}
\]
The original host bounds give $\operatorname{rd}(E)^2\ge2f$
by AH16.  Also $e\ge2$, $2^e\mid\operatorname{ord}$, and
$f=2\operatorname{ord}$, so $f\ge8$.  Positivity yields
$\operatorname{rd}(E)\ge4$, hence $v\ge2$.
Therefore $4k^4v^2\le(kv)^4$.  Using $d_\tau\ge0$,
$T\ge2$, and GLS4 in (GLS.47) gives exactly
\[
 \boxed{\quad
 \mu_{\max}(A)\le
   \sum_\tau\left(4d_\tau\log(kv)
          +k^2\log(h_\tau+k)+|s_\tau|\log T\right).
 \quad}                                                       \tag{GLS.48}
\]
This is \texttt{SlopeMax D u 4}, on every original $D,u$ and
every given \texttt{ColVal D u}.  It uses no Hida size guard.
The transference input in GLS8 is an explicit, fully matched
published theorem; all subsequent arithmetic, ideal, and metric
comparisons have been proved here.
% END GLS11

\subsection{Uniform installation and the remaining final assembly}

% BEGIN GLS12
\paragraph{GLS12: all three literal GL4 family entries.}
\label{claim:GLS12}
For every original family $\mathrm{Fm}:\texttt{FamP}\,\varepsilon$,
GLS7 and GLS11 apply to every $D=\mathrm{Fm}.D(x,i)$ and every
$u,\mathrm{CV}$.  Threshold zero consequently gives
\[
 \texttt{FamilySlopeLowerOpenPart}\,\mathrm{Fm},\qquad
 \texttt{FamilySlopeMaxOpenPart}\,\mathrm{Fm}\,4.
 \tag{GLS.49}
\]
Their \texttt{fam0} quantifiers are
$\exists x_0\ \forall x\ge x_0\ \forall (h:\mathrm{Fm}.x_1<x)
\ \forall i$, followed by the pointwise predicate.  The argument
works for all such $x,i$, so these are universal \texttt{Closed}
statements on any family class, with no change of $\varepsilon$,
$C_H$, or the original member.
Explicitly, \texttt{Closed} means
\[
 \exists\varepsilon_1>0\ \forall\varepsilon\in(0,\varepsilon_1]
 \ \exists\mathcal K\ \forall\mathrm{Fm},C_H,\quad
 \mathcal A(\mathrm{Fm})\wedge\mathrm{Fm}.\mathrm{Large}(\mathcal K,C_H)
       \Longrightarrow P(\mathrm{Fm},C_H).
\]
For each assertion in (GLS.49) one may take
$\varepsilon_1=1$ and the original all-zero \texttt{TransportedAlt}
tuple $\mathcal K$; the hypotheses of \texttt{Large} are unused.

Suppose the same family has \texttt{RealPlaceSecondLayerConditions}
at threshold $a$.  Thus, for every $x\ge a$ and every $i$, it
supplies the original real-place tie and one
$\mathcal F_{x,i}:\texttt{ColsTiedSym}(\mathrm{Fm}.D(x,i))$,
whose forms have their supplied finite characters for every word.
RSL7 applied to each original $\mathcal F_{x,i}.\mathrm{forms}(u)$
already gives \texttt{admS} pointwise, for every $u$; their conjunction
is \texttt{adm3}.  Thus one may take $b=a$.  The source's
\texttt{family\_conditions\_of\_layer\_two} also permits a separate
eventual stabilizer threshold $b$ with the guard
$\mathrm{Fm}.x_1<x$; the following common choice covers both routes:
\[
                    a_{\rm CV}=\max\{a,b,\mathrm{Fm}.x_1+1\}.
 \tag{GLS.50}
\]
For $x\ge a_{\rm CV}$, (GLS.15) applied to this very
$\mathcal F_{x,i}$ gives
\[
 \forall i:\mathrm{Fm}.Ix(x)\ \forall u,\quad
            \texttt{Nonempty}(\texttt{ColVal}(\mathrm{Fm}.D(x,i),u)).
 \tag{GLS.51}
\]
This is the literal unguarded \texttt{tCV} field: its threshold
is chosen before $i,u$, and (GLS.50) explicitly removes the strict
$x_1$ guard.  Any \texttt{Closed} second-layer input therefore
gives \texttt{Closed tCV} on the same class.

Equations (GLS.49)--(GLS.51) close three stated implication tasks:
the lower slope, maximum slope, and original column-value dictionary.
Their audit identifiers are respectively
\texttt{TASK\_GL4\_SLOPE\_LOWER},
\texttt{TASK\_GL4\_SLOPE\_MAX}, and
\texttt{TASK\_GL4\_COLUMN\_VALUES}.
The last retains the original same-selector second-layer producer.
The two slope inequalities are conditional on their given
\texttt{ColVal}, as their source predicates require.
% END GLS12

% BEGIN GLS13
\paragraph{GLS13: the complete eleven-field packet and its precise frontier.}
\label{claim:GLS13}
Let $\mathcal A$ be a family class and retain
\[
 \forall\mathrm{Fm}\in\mathcal A,\quad\texttt{Duty}(\mathrm{Fm}),
 \qquad
 \texttt{Closed}_{D_3(3,\mathcal A)}
                   (\texttt{RealPlaceSecondLayerConditions}).
 \tag{GLS.52}
\]
All the following applications concern that same class, family,
and second-layer witness.  MF19 supplies \texttt{h2b};
RSL1--RSL3, using the stated mass and volume inputs, supply
\texttt{h2v}; DSZ1--DSZ3 supply \texttt{h3}.
GLS12 supplies \texttt{hZa}, \texttt{hZb} with $C_5=4$, and
\texttt{tCV}.  The original finite character from (GLS.52) and
RSL7--RSL10 supply \texttt{hW} and \texttt{h}$\mu$.
Finally GLE5, GLE9, and GLE6 supply \texttt{hI}, \texttt{hR}, and
\texttt{hB}, respectively, retaining the original \texttt{FamU}
guard \texttt{HidaSetting}.  GLE12 proves their uniform installation.
These are exactly the eleven fields, in order, of
\[
                    \texttt{FamGLmin}(D_3(3,\mathcal A),4).
 \tag{GLS.53}
\]
Each eventual intersection uses a maximum of finitely many existing
thresholds.  The maximum is taken before the member $i$, word $u$,
coefficient place, or level place is quantified.  There is no need
to choose a second form record or change any field of the family.
When combining the \texttt{Closed} assertions, take the minimum
of their positive $\varepsilon_1$ bounds and the componentwise
supremum of their finitely many \texttt{TransportedAlt} tuples.
The original monotonicity of \texttt{Large} recovers each input
for the very same $\mathrm{Fm},C_H$.  These choices precede the
family and its members, as the displayed quantifier requires.
Notice also that (GLS.53), like the source's guarded fields, does
not itself require \texttt{TuSize}.  That condition supplies
\texttt{SettingEventually} when those guarded fields are consumed.

The smaller fullness/rank conclusion remains the already proved
RSL11--RSL12 and GLE route; it does not depend on GLS4--GLS11.
Membership in $D_3(3,\mathcal A)$ supplies the same family's
\texttt{TuSize}.  The Duty clause supplies that family's
\texttt{ChosenWall}.  These are exactly the two projections used
in the source's \texttt{link\_gl} after constructing (GLS.53).

We keep the final analytic assembly distinct.  The source's
\texttt{linkGLOBAL7\_min} first transforms (GLS.53) into
\texttt{PartFiveFamilyBounds}, \texttt{ConclusionF5},
\texttt{ConclusionF6}, \texttt{family\_property\_and}, and
\texttt{CongruenceLatticeLengthBound}.  Its
\texttt{linkGLOBAL7\_items} then combines the original count,
two slope families, rational ranks, lattices, frames, loss windows,
characters, and size conditions in \texttt{linkCore}.
Giving the names of those source maps is not a human proof of all
their numerical estimates and input identifications.  Accordingly
the remaining B obligation here is the actual numerical/identity
installation from this eleven-field packet to
\[
 \texttt{TuSize}\to\texttt{DGamSize}\to
                    \texttt{LinkGLOBAL}(\mathrm{Fm},2,C_H,3,2,0),
 \tag{GLS.54}
\]
through \texttt{linkGLOBAL7\_min},
\texttt{linkGLOBAL7\_items}, and \texttt{FamGLmin.link8}.
The final \texttt{link8} projection retains (GLS.52)'s
\texttt{ChosenWall}, not an independently chosen covering family.
The coefficient $C_5=4$ in (GLS.53), $C_{\rm abs}=2$, $C_0=3$,
$e_D=2$, and $e_L=0$ must remain the same in that task.
The common compatible selector and (GLS.52) remain upstream
producer obligations; neither is asserted to exist by this chapter.
% END GLS13
